\message{ !name(regular-and-exact-categories.tex)}\documentclass{article}
%
\newcommand{\cat}[1]{\mathsf{#1}}
\newcommand{\univ}[1]{\mathcal{#1}}
\newcommand{\path}{\equiv}
\newcommand{\id}[1]{\mathsf{#1}}
\newcommand{\realizes}[3]{#1 \Vdash_{#3} #2}
\usepackage{amsmath}
\usepackage{amsthm}
\usepackage{amssymb}
\usepackage{tikz}
\usepackage{quiver}
\newtheorem{theorem}{Theorem}
\newtheorem{prop}[theorem]{Proposi{t}ion}
\newtheorem{lemma}[theorem]{Lemma}
\theoremstyle{plain}
\newtheorem{corollary}[theorem]{Corollary}
\theoremstyle{remark}
\newtheorem{defn}[theorem]{Def{i}ni{t}ion}
\usepackage{xcolor}
\usepackage{natbib}

% From CAN package "orcidlink"
\usepackage{scalerel}
\usepackage{tikz}
\usetikzlibrary{svg.path}
\definecolor{orcidlogocol}{HTML}{A6CE39}
\tikzset{
  orcidlogo/.pic={
    \fill[orcidlogocol] svg{M256,128c0,70.7-57.3,128-128,128C57.3,256,0,198.7,0,128C0,57.3,57.3,0,128,0C198.7,0,256,57.3,256,128z};
    \fill[white] svg{M86.3,186.2H70.9V79.1h15.4v48.4V186.2z}
                 svg{M108.9,79.1h41.6c39.6,0,57,28.3,57,53.6c0,27.5-21.5,53.6-56.8,53.6h-41.8V79.1z M124.3,172.4h24.5c34.9,0,42.9-26.5,42.9-39.7c0-21.5-13.7-39.7-43.7-39.7h-23.7V172.4z}
                 svg{M88.7,56.8c0,5.5-4.5,10.1-10.1,10.1c-5.6,0-10.1-4.6-10.1-10.1c0-5.6,4.5-10.1,10.1-10.1C84.2,46.7,88.7,51.3,88.7,56.8z};
  }
}

\newcommand\orcidicon[1]{\href{https://orcid.org/#1}{\mbox{\scalerel*{
\begin{tikzpicture}[yscale=-1,transform shape]
\pic{orcidlogo};
\end{tikzpicture}
}{|}}}}
% end orcidlink
\usepackage{hyperref}
% notation 
% end notationa
\begin{document}

\message{ !name(regular-and-exact-categories.tex) !offset(116) }
\section{$\mathsf{Asm}$ is regular}
\label{sec:asm-regular}

Let us proceed to prove that the category of assemblies is regular. Throughout this section, assume that $A : \mathcal{U}$ is the underlying set of an arbitrary total combinatory algebra and $\mathcal{V} \; \mathcal{W}$ are arbitrary universes.

\begin{defn}
  An \textbf{assembly} over a set $X : \mathcal{V}$ consists of a relation $\Vdash : A \to X \to \mathcal{W}$ that is proposition valued along with a witness of the surjectivity of $\Vdash$ so that,
  $\forall (x : X) \; \exists (a : A) \; \realizes{a}{x}{X}$
  \end{enumerate}
\end{defn}

Immediately we see that the space of all assemblies is a set.

Morphisms between assemblies are defined as follows.

\begin{defn}
  An \textbf{assembly morphism} $X \to Y$ is a map $f$ from the underlying set of $X$ to the underlying set of $Y$ that is \textbf{tracked}, we should merely have a tracker $r$ such that
  $\forall (x : X) (x_r : A) \to \realizes{x_r}{x}{X} \to \realizes{r \; x_r}{f(x)}{Y}$.
\end{defn}

It is useful to have a characterisation of the path space of assembly morphisms. The following
lemma is trivial but incredibly useful.

\begin{lemma}
  Given assemblies $X$ and $Y$ and assembly morphisms $f \; g : X \to Y$ such that $\forall (x : X)\; f(x) = g(x)$ then $f = g$.
\end{lemma}
\begin{proof}
  Trivial by the definition of assembly morphism and function extensionality.
\end{proof}

\bibliographystyle{plainnat}
\bibliography{realisability}

\message{ !name(regular-and-exact-categories.tex) !offset(114) }

\end{document}
